\documentclass{amsart}
% For dealing with references we use the comment environment
\usepackage{verbatim}
\newenvironment{reference}{\comment}{\endcomment}
%\newenvironment{reference}{}{}
\newenvironment{slogan}{\comment}{\endcomment}
\newenvironment{history}{\comment}{\endcomment}

% For commutative diagrams we use Xy-pic
\usepackage[all]{xy}

% We use 2cell for 2-commutative diagrams.
\xyoption{2cell}
\UseAllTwocells

% We use multicol for the list of chapters between chapters
\usepackage{multicol}

% This is generally recommended for better output
\usepackage{lmodern}
\usepackage[T1]{fontenc}

% For cross-file-references

% Package for hypertext links:
\usepackage{hyperref}

% For any local file, say "hello.tex" you want to link to please

% Theorem environments.
%
\theoremstyle{plain}
\newtheorem{theorem}[subsection]{Theorem}
\newtheorem{proposition}[subsection]{Proposition}
\newtheorem{lemma}[subsection]{Lemma}

\theoremstyle{definition}
\newtheorem{definition}[subsection]{Definition}
\newtheorem{example}[subsection]{Example}
\newtheorem{exercise}[subsection]{Exercise}
\newtheorem{situation}[subsection]{Situation}

\theoremstyle{remark}
\newtheorem{remark}[subsection]{Remark}
\newtheorem{remarks}[subsection]{Remarks}

\numberwithin{equation}{subsection}

% Macros
%
\def\lim{\mathop{\mathrm{lim}}\nolimits}
\def\colim{\mathop{\mathrm{colim}}\nolimits}
\def\Spec{\mathop{\mathrm{Spec}}}
\def\Hom{\mathop{\mathrm{Hom}}\nolimits}
\def\Ext{\mathop{\mathrm{Ext}}\nolimits}
\def\SheafHom{\mathop{\mathcal{H}\!\mathit{om}}\nolimits}
\def\SheafExt{\mathop{\mathcal{E}\!\mathit{xt}}\nolimits}
\def\Sch{\mathit{Sch}}
\def\Mor{\mathop{\mathrm{Mor}}\nolimits}
\def\Ob{\mathop{\mathrm{Ob}}\nolimits}
\def\Sh{\mathop{\mathit{Sh}}\nolimits}
\def\NL{\mathop{N\!L}\nolimits}
\def\CH{\mathop{\mathrm{CH}}\nolimits}
\def\proetale{{pro\text{-}\acute{e}tale}}
\def\etale{{\acute{e}tale}}
\def\QCoh{\mathit{QCoh}}
\def\Ker{\mathop{\mathrm{Ker}}}
\def\Im{\mathop{\mathrm{Im}}}
\def\Coker{\mathop{\mathrm{Coker}}}
\def\Coim{\mathop{\mathrm{Coim}}}

% Boxtimes
%
\DeclareMathSymbol{\boxtimes}{\mathbin}{AMSa}{"02}

%
% Macros for moduli stacks/spaces
%
\def\QCohstack{\mathcal{QC}\!\mathit{oh}}
\def\Cohstack{\mathcal{C}\!\mathit{oh}}
\def\Spacesstack{\mathcal{S}\!\mathit{paces}}
\def\Quotfunctor{\mathrm{Quot}}
\def\Hilbfunctor{\mathrm{Hilb}}
\def\Curvesstack{\mathcal{C}\!\mathit{urves}}
\def\Polarizedstack{\mathcal{P}\!\mathit{olarized}}
\def\Complexesstack{\mathcal{C}\!\mathit{omplexes}}
% \Pic is the operator that assigns to X its picard group, usage \Pic(X)
% \Picardstack_{X/B} denotes the Picard stack of X over B
% \Picardfunctor_{X/B} denotes the Picard functor of X over B
\def\Pic{\mathop{\mathrm{Pic}}\nolimits}
\def\Picardstack{\mathcal{P}\!\mathit{ic}}
\def\Picardfunctor{\mathrm{Pic}}
\def\Deformationcategory{\mathcal{D}\!\mathit{ef}}

\setlength{\emergencystretch}{3em}
\begin{document}
\title[Stacks Project, Tag 0F5U]{1383 --- Gabriel-Popescu exact localization of every Grothendieck abelian category by a module category}
\maketitle
\noindent Copyright (C) 2005--2025 Johan de Jong. Stacks Project authors. Source: \href{https://stacks.math.columbia.edu/tag/0F5U}{Tag 0F5U}.

\noindent Editorial difficulty note (not author text). Difficulty H2: The author realizes the category through modules over the endomorphism ring of a possibly nonprojective generator. The main difficulty is proving the counit injective from finite-generator tests and then proving exactness of the left adjoint using AB5 and finite free submodules; this is a substantive structural theorem beyond qualification category or homological algebra..

\noindent Editorial provenance note (not author text). History: excerpt from revision \texttt{a0444\allowbreak 6e57e\allowbreak c1fbc\allowbreak 252a8\allowbreak 71afc\allowbreak ec775\allowbreak 2fb28\allowbreak 07b14}, injectives.tex, lines 2660--2702. Reference presentation was adapted; original-author context and core support blocks were retained or added. The mathematical wording is unchanged. Licensed under GNU FDL 1.2 or later; no invariant sections or cover texts. See ../licenses/COPYING. Context blocks reproduce author definitions and supporting results; they are not additional counted problems.

\begin{definition}
\label{definition-grothendieck-conditions}
Let $\mathcal{A}$ be an abelian category. We name some conditions
\begin{enumerate}
\item[AB3] $\mathcal{A}$ has direct sums,
\item[AB4] $\mathcal{A}$ has AB3 and direct sums are exact,
\item[AB5] $\mathcal{A}$ has AB3 and filtered colimits are exact.
\end{enumerate}
Here are the dual notions
\begin{enumerate}
\item[AB3*] $\mathcal{A}$ has products,
\item[AB4*] $\mathcal{A}$ has AB3* and products are exact,
\item[AB5*] $\mathcal{A}$ has AB3* and cofiltered limits are exact.
\end{enumerate}
We say an object $U$ of $\mathcal{A}$ is a {\it generator} if
for every $N \subset M$, $N \not = M$ in $\mathcal{A}$ there exists a morphism
$U \to M$ which does not factor through $N$.
We say $\mathcal{A}$ is a {\it Grothendieck abelian category} if
it has AB5 and a generator.
\end{definition}

\begin{definition}
\label{definition-size}
Let $\mathcal{A}$ be a Grothendieck abelian category.
Let $M$ be an object of $\mathcal{A}$.
The {\it size} $|M|$ of $M$ is the cardinality of the set of subobjects
of $M$.
\end{definition}

\begin{theorem}[Adjoint functor theorem]
\label{categories-theorem-adjoint-functor}
Let $G : \mathcal{C} \to \mathcal{D}$ be a functor of big categories.
Assume $\mathcal{C}$ has limits, $G$ commutes with them, and for
every object $y$ of $\mathcal{D}$ there exists a set of pairs
$(x_i, f_i)_{i \in I}$ with $x_i \in \Ob(\mathcal{C})$,
$f_i \in \Mor_\mathcal{D}(y, G(x_i))$ such that for any
pair $(x, f)$ with $x \in \Ob(\mathcal{C})$,
$f \in \Mor_\mathcal{D}(y, G(x))$ there are an $i$ and a morphism
$h : x_i \to x$ such that $f = G(h) \circ f_i$.
Then $G$ has a left adjoint $F$.
\end{theorem}

\begin{lemma}
\label{lemma-F-G-monos}
Let $f : M \to G(A)$ be an injective map in $\text{Mod}_R$.
Then the adjoint map $f' : F(M) \to A$ is injective too.
\end{lemma}

\begin{proof}
Choose a map $R^{\oplus n} \to M$ and consider the corresponding map
$U^{\oplus n} \to F(M)$. Consider a map $v : U \to U^{\oplus n}$
such that the composition $U \to U^{\oplus n} \to F(M) \to A$ is $0$.
Then this arrow $v : U \to U^{\oplus n}$ is an element
$v$ of $R^{\oplus n}$ mapping to zero in $G(A)$. Since $f$ is injective,
we conclude that $v$ maps to zero in $M$ which means that
$U \to U^{\oplus n} \to F(M)$ is zero by construction of $F(M)$
in the proof of Lemma \ref{lemma-gabriel-popescu-left-adjoint}.
Since $U$ is a generator we conclude that
$$
\Ker(U^{\oplus n} \to F(M) \to A) = \Ker(U^{\oplus n} \to F(M))
$$
To finish the proof we choose a surjection $\bigoplus_{i \in I} R \to M$
and we consider the corresponding surjection
$$
\pi : \bigoplus\nolimits_{i \in I} U \longrightarrow F(M)
$$
To prove $f'$ is injective it suffices to show that
$\Ker(\pi) = \Ker(f' \circ \pi)$ as subobjects of $\bigoplus_{i \in I} U$.
However, now we can write $\bigoplus_{i \in I} U$ as the filtered colimit
of its subobjects $\bigoplus_{i \in I'} U$ where $I' \subset I$
ranges over the finite subsets. Since filtered colimits are
exact by AB5 for $\mathcal{A}$, we see that
$$
\Ker(\pi) =
\colim_{I' \subset I\text{ finite}}
\left(\bigoplus\nolimits_{i \in I'} U\right)
\bigcap \Ker(\pi)
$$
and
$$
\Ker(f' \circ \pi) =
\colim_{I' \subset I\text{ finite}}
\left(\bigoplus\nolimits_{i \in I'} U\right)
\bigcap \Ker(f' \circ \pi)
$$
and we get equality because the same is true for each $I'$ by
the first displayed equality above.
\end{proof}


\begin{lemma}
\label{lemma-set-iso-classes-bounded-size}
Let $\mathcal{A}$ be a Grothendieck abelian category with generator $U$.
\begin{enumerate}
\item If $|M| \leq \kappa$, then $M$ is the quotient of a direct
sum of at most $\kappa$ copies of $U$.
\item For every cardinal $\kappa$ the isomorphism classes
of objects $M$ with $|M| \leq \kappa$ form a set.
\end{enumerate}
\end{lemma}

\begin{proof}
For (1) choose for every proper subobject $M' \subset M$ a morphism
$\varphi_{M'} : U \to M$ whose image is not contained in $M'$. Then
$\bigoplus_{M' \subset M} \varphi_{M'} : \bigoplus_{M' \subset M} U \to M$
is surjective. It is clear that (1) implies (2).
\end{proof}


\begin{lemma}
\label{lemma-set-of-subobjects}
Let $\mathcal{A}$ be an abelian category with a generator $U$ and
$X$ an object of $\mathcal{A}$. If $\kappa$ is the cardinality of
$\Mor(U, X)$ then
\begin{enumerate}
\item There does not exist a strictly increasing
(or strictly decreasing) chain of subobjects
of $X$ indexed by a cardinal bigger than $\kappa$.
\item If $\alpha$ is an ordinal of cofinality $> \kappa$
then any increasing (or decreasing) sequence of subobjects
of $X$ indexed by $\alpha$ is eventually constant.
\item The cardinality of the set of subobjects of $X$
is $\leq 2^\kappa$.
\end{enumerate}
\end{lemma}

\begin{proof}
For (1) assume $\kappa' > \kappa$ is a cardinal and assume
$X_i$, $i \in \kappa'$ is strictly increasing. Then take for
each $i$ a $\phi_i \in \Mor(U, X)$ such that $\phi_i$ factors through
$X_{i + 1}$ but not through $X_i$. Then the morphisms $\phi_i$
are distinct, which contradicts the definition of $\kappa$.

\medskip\noindent
Part (2) follows from the definition of cofinality and (1).

\medskip\noindent
Proof of (3). For any subobject $Y \subset X$
define $S_Y \in \mathcal{P}(\Mor(U, X))$ (power set) as
$S_Y = \{\phi \in \Mor(U,X) : \phi)\text{ factors through }Y\}$.
Then $Y = Y'$ if and only if $S_Y = S_{Y'}$. Hence the cardinality
of the set of subobjects is at most the cardinality of this power set.
\end{proof}


\medskip\noindent
Let $\mathcal{A}$ be a Grothendieck abelian category and let $U$ be a
generator for $\mathcal{A}$, see
Definition \ref{definition-grothendieck-conditions}.
Let $R = \Hom_\mathcal{A}(U, U)$. Consider the functor
$G : \mathcal{A} \to \text{Mod}_R$ given by
$$
G(A) = \Hom_\mathcal{A}(U, A)
$$
endowed with its canonical right $R$-module structure.

\begin{lemma}
\label{lemma-gabriel-popescu-left-adjoint}
The functor $G$ above has a left adjoint
$F : \text{Mod}_R \to \mathcal{A}$.
\end{lemma}

\begin{proof}
We will give two proofs of this lemma.

\medskip\noindent
The first proof will use the adjoint functor theorem, see
Categories, Theorem \ref{categories-theorem-adjoint-functor}.
Observe that $G : \mathcal{A} \to \text{Mod}_R$ is left exact and sends
products to products. Hence $G$ commutes with limits. To check the set
theoretical condition in the theorem, suppose that $M$ is an object of
$\text{Mod}_R$. Choose a suitably large cardinal $\kappa$ and denote $E$
a set of objects of $\mathcal{A}$ such that every object $A$ with
$|A| \leq \kappa$ is isomorphic to an element of $E$. This is possible
by Lemma \ref{lemma-set-iso-classes-bounded-size}. Set
$I = \coprod_{A \in E} \Hom_R(M, G(A))$.
We think of an element $i \in I$ as a pair $(A_i, f_i)$.
Finally, let $A$ be an arbitrary object of $\mathcal{A}$
and $f : M \to G(A)$ arbitrary. We are going to think of
elements of $\Im(f) \subset G(A) = \Hom_\mathcal{A}(U, A)$
as maps $u : U \to A$. Set
$$
A' = \Im(\bigoplus\nolimits_{u \in \Im(f)} U \xrightarrow{u} A)
$$
Since $G$ is left exact, we see that $G(A') \subset G(A)$
contains $\Im(f)$ and we get $f' : M \to G(A')$ factoring $f$.
On the other hand, the object $A'$ is
the quotient of a direct sum of at most $|M|$ copies of $U$.
Hence if $\kappa = |\bigoplus_{|M|} U|$, then we see that $(A', f')$
is isomorphic to an element $(A_i, f_i)$ of $E$ and we conclude that $f$
factors as $M \xrightarrow{f_i} G(A_i) \to G(A)$ as desired.

\medskip\noindent
The second proof will give a construction of $F$ which will show
that ``$F(M) = M \otimes_R U$'' in some sense. Namely, for any
$R$-module $M$ we can choose a resolution
$$
\bigoplus\nolimits_{j \in J} R \to
\bigoplus\nolimits_{i \in I} R \to
M \to 0
$$
Then we define $F(M)$ by the corresponding exact sequence
$$
\bigoplus\nolimits_{j \in J} U \to
\bigoplus\nolimits_{i \in I} U \to
F(M) \to 0
$$
This construction is independent of the choice of the resolution
and is functorial; we omit the details.
For any $A$ in $\mathcal{A}$ we obtain an exact sequence
$$
0 \to \Hom_\mathcal{A}(F(M), A) \to
\prod\nolimits_{i \in I} G(A) \to
\prod\nolimits_{j \in J} G(A)
$$
which is isomorphic to the sequence
$$
0 \to \Hom_R(M, G(A)) \to
\Hom_R(\bigoplus\nolimits_{i \in I} R, G(A)) \to
\Hom_R(\bigoplus\nolimits_{j \in J} R, G(A))
$$
which shows that $F$ is the left adjoint to $G$.
\end{proof}


\begin{theorem}
\label{theorem-gabriel-popescu}
Let $\mathcal{A}$ be a Grothendieck abelian category. Then there exists
a (noncommutative) ring $R$ and functors $G : \mathcal{A} \to \text{Mod}_R$
and $F : \text{Mod}_R \to \mathcal{A}$ such that
\begin{enumerate}
\item $F$ is the left adjoint to $G$,
\item $G$ is fully faithful, and
\item $F$ is exact.
\end{enumerate}
Moreover, the functors are the ones constructed above.
\end{theorem}

\begin{proof}
We first prove $G$ is fully faithful, or equivalently that
$F \circ G \to \text{id}$ is an isomorphism, see
Categories, Lemma \href{https://stacks.math.columbia.edu/tag/07RB}{lemma adjoint fully faithful (Tag 07RB)}.
First, given an object $A$ the map $F(G(A)) \to A$ is surjective,
because every map of $U \to A$ factors through $F(G(A))$ by construction.
On the other hand, the map $F(G(A)) \to A$ is the adjoint of the
map $\text{id} : G(A) \to G(A)$ and hence injective by
Lemma \ref{lemma-F-G-monos}.

\medskip\noindent
The functor $F$ is right exact as it is a left adjoint.
Since $\text{Mod}_R$ has enough projectives, to show that
$F$ is exact, it is enough to show that the first left derived
functor $L_1F$ is zero. To prove $L_1F(M) = 0$ for some $R$-module $M$
choose an exact sequence $0 \to K \to P \to M \to 0$
of $R$-modules with $P$ free. It suffices to show $F(K) \to F(P)$
is injective. Now we can write this sequence as a filtered
colimit of sequences $0 \to K_i \to P_i \to M_i \to 0$
with $P_i$ a finite free $R$-module: just write $P$ in this
manner and set $K_i = K \cap P_i$ and $M_i = \Im(P_i \to M)$.
Because $F$ is a left adjoint it commutes
with colimits and because $\mathcal{A}$ is a Grothendieck
abelian category, we find that $F(K) \to F(P)$
is injective if each $F(K_i) \to F(P_i)$ is injective.
Thus it suffices to check $F(K) \to F(P)$
is injective when $K \subset P = R^{\oplus n}$.
Thus $F(K) \to U^{\oplus n}$ is injective by an application
of Lemma \ref{lemma-F-G-monos}.
\end{proof}
\end{document}
